\section{引言}\label{sec:intro}

非阿贝尔规范理论所描述的物理可以通过多种方式加以研究。依据具体情形
与所要回答的问题，所需信息可以从定域场的关联函数、Wilson 圈的期望值
或者薛定谔泛函等对象中提取。

正如文献~\cite{WilsonFlow,Villasimius} 所解释的，梯度流为进一步探究这些理论
提供了新的途径，并使人们能够理解它们某些其他方式难以把握的性质
（在格点规范理论中，梯度流被称为 Wilson 流，
而在数学文献中它通常被称作 Yang--Mills 梯度流）。
显然，有物理意义的探测量必须不受紫外发散的危害，
或者说必须能通过一个良定义的重正化手续将这些发散消去。

就梯度流而言，已有一些证据表明由流生成的规范场不需要重正化
\cite{WilsonFlow,Villasimius}，但到目前为止，
尚未有人给出正流时间处不存在紫外发散的正式证明。
本文的目的正是通过对流的微扰论作全阶分析来填补这一空白。
与 Zinn--Justin 和 Zwanziger 讨论的密切相关的 Langevin 方程的重正化
\cite{Zinn,ZinnZwanziger} 相比，这里有两点重要差别：
其一是流方程中没有噪声项，其二是规范场的初始分布
（即由所考虑理论的泛函积分给出）并没有被忽略。

本文给出的微扰分析适用于具有任意紧致单群规范群和任意物质多重态的
可重正化规范理论。不过，为了使讨论尽可能地简单，
我们只考虑采用维数正规化的纯 $\SUn$ 规范理论，
向其他情形的推广是直接的（见第 \ref{sec:misc} 节）。

在接下来的三节中，我们将导出由梯度流生成的规范场关联函数的费曼规则，
并证明它们就是一个带额外维度（流时间）的定域场论的费曼规则。
随后可以借助幂次计数与 BRS 对称性确立正流时间处关联函数的有限性
（第 \ref{sec:brs}、\ref{sec:fin} 节），
但作为示例，我们先在第 \ref{sec:sample} 节中算出一组单圈图的发散部分。
